\documentclass{article}
\usepackage{amsmath,amssymb,xurl}
\usepackage[hidelinks]{hyperref}
% Shared commands for notes in docs/paper-gaps/.

\DeclareUnicodeCharacter{2080}{\ifmmode{}_{0}\else\textsubscript{0}\fi}
\DeclareUnicodeCharacter{2081}{\ifmmode{}_{1}\else\textsubscript{1}\fi}
\DeclareUnicodeCharacter{2082}{\ifmmode{}_{2}\else\textsubscript{2}\fi}
\DeclareUnicodeCharacter{2099}{\ifmmode{}_{n}\else\textsubscript{n}\fi}

\newcommand{\Lean}{\textsc{Lean}}
\ExplSyntaxOn
\NewDocumentCommand{\leanid}{m}{
  \ifmmode
    \textsf{\detokenize{#1}}
  \else
    \group_begin:
    \urlstyle{sf}
    \tl_set:Nn \l_tmpa_tl {#1}
    \tl_replace_all:Nnn \l_tmpa_tl {\_} {_}
    \exp_args:NV \path \l_tmpa_tl
    \group_end:
  \fi
}
\ExplSyntaxOff
\providecommand{\tr}{\operatorname{tr}}
\newcommand{\tnleanIssueBase}{https://github.com/LionSR/TNLean/issues/}
\newcommand{\tnleanPrBase}{https://github.com/LionSR/TNLean/pull/}
\newcommand{\tnleanDocBase}{https://sirui-lu.com/TNLean/docs/}
\newcommand{\ghissue}[1]{\href{\tnleanIssueBase#1}{GitHub issue \##1}}
\newcommand{\ghpr}[1]{\href{\tnleanPrBase#1}{PR \##1}}
\newcommand{\doclink}[1]{\href{\tnleanDocBase#1.html}{\path{#1}}}

% Dirac notation and the identity operator, as in blueprint/src/macros/common.tex.
% \providecommand keeps note-local definitions authoritative.
\usepackage{dsfont}
\providecommand{\ket}[1]{|#1\rangle}
\providecommand{\bra}[1]{\langle #1|}
\providecommand{\braket}[2]{\langle #1 | #2 \rangle}
\providecommand{\ketbra}[2]{|#1\rangle\!\langle #2|}
\providecommand{\Id}{\mathds{1}}
\providecommand{\one}{\mathds{1}}
\providecommand{\C}{\mathbb{C}}
\providecommand{\MN}[1]{M_{#1}(\C)}
\providecommand{\MD}{M_D(\C)}

% Verdict marker: \gapnote{<kind>}{<status>}. Kind is the policy
% classification (clarification, local-correction, scope-restriction,
% unfaithful, false-source, open-gap); status is open, wip (elimination
% underway), resolved, or historical. Machine-read by the site index;
% prints a verdict line.
\newcommand{\gapnote}[2]{\par\noindent\textbf{Verdict:} #1 (#2).\par}

\title{Noncentral Reduced Density from an Untwisted Boundary Component}
\author{}
\date{}
\begin{document}
\maketitle
\gapnote{scope-restriction}{open}

Schuch, Cirac, and P\'erez-Garc\'ia,
\href{https://arxiv.org/abs/1001.3807v3}{arXiv:1001.3807v3},
Lemma 6.14 and Theorem 6.9, source
\path{Papers/1001.3807/paper_v3.tex}, lines 1935--1990 and 2199--2214,
motivate comparing an open flux pair with the vacuum away from its ends.
For noncentral insertions, identity on the observed incident bonds alone
is insufficient on arbitrary graphs. The complementary boundary endpoints
must also be connected by an appropriate untwisted component.

The theorem
\leanid{TNLean.PEPS.regularPhysicalCutMatrix_normalized_reduced_eq_of_untwisted_boundary_component}
proves the actual normalized reduced-density equality under the following
explicit structural assumptions: a finite graph and finite group, regular
group-isometric physical sites, connected observed and complementary
induced graphs, nonempty boundary, identity insertions on every observed
incident bond, and an untwisted complementary component containing every
complementary boundary endpoint. The two actual cuts are required nonzero.
There is no centrality, supplied spanning tree, global gauge, contraction
factorization, Gram identity, or density equality hypothesis.

The anchored tree construction supplies a global gauge fixing the observed
region and all its incident operators while preserving the actual closed
vector. Its complementary tree then supplies the connected common
untwisted graph needed in the projected mixed Gram calculation. Applying
that calculation to the transformed insertion and the vacuum gives a common
normalized physical density. Positive scalar normalizations are converted
to trace normalization of the actual coefficient matrices. Unnormalized
reduced matrices are not asserted equal; noncentral insertions may change
their scalar factors.

This is a conditional noncentral contraction theorem. For a native torus
open flux pair, the complementary boundary component must still be derived
from the actual exterior collar and the chosen dual detour. The native
regular-isometric nonvanishing results can then discharge the explicit
nonzero-cut assumptions. The theorem does not infer the component from
endpoint exclusion or equate arbitrary endpoint-free annular observations
with vacuum. These remaining geometric source obligations stay under
\ghissue{8703}; the full-lattice Hamiltonian scope of \ghissue{8947} is
unchanged.
\end{document}
